\scanpage{21}

tels termes. Donc $c_{ij}\equiv0\ (p)$ dans ce cas, ce qui prouve (remplaçant $i,j$ par $u,v$) que $c_{u,v}\equiv0\ (p)$ donc encore $c_{ij}\equiv c_{u,v}c_{r,s}\ (p)$.

Le lemme 1 résulte du \unsure{corollaire} par une récurrence immédiate sur la longueur $e$ des développements $p$-adiques écrits de $u$ et de $v$.

\uline{Cor. 2} Soit $N$ un entier \inserted{$\geq0$} tel que $N+1\equiv0\ (p)$, posons
\figure{\begin{tikzpicture}\draw (0,0)--(0,-2.7);\end{tikzpicture}}
\[
N+1=(N'+1)p\quad\text{donc}\quad N=(N'+1)p-1=N'p+(p-1)
\]
(de sorte que à décalage de 1 près, les $N+1$ et $N'+1$ ont $=$ chiff. du développement $p$-adique, ainsi que $N$ et $N'$ (quand on fait abstraction du premier chiff. p.a. de $N$). Soit donc $\varepsilon\geq1$, on a
\[
\begin{aligned}
c_{N+1,p^\varepsilon-1}\deleted{\unsure{\not\equiv0\ (p)\Longleftrightarrow}}\inserted{\equiv}
c_{N'+1,p^{\varepsilon-1}-1}\deleted{\unsure{\not\equiv0\ (p)}}\deleted{\unsure{\ne0}}
\left\{\begin{array}{ll}
\equiv1\ (p)&\text{si }N'+1\equiv0\ (p^{\varepsilon-1})\\
\equiv0\ (p)&\text{si }N'+1\not\equiv0\ (p^{\varepsilon-1})
\end{array}\right.\\
c_{N,p^\varepsilon}\deleted{\unsure{\not\equiv0\ (p)\Longleftrightarrow}}\inserted{\equiv}
c_{N',p^{\varepsilon-1}}\deleted{\unsure{\not\equiv0}}\ (p)\deleted{\unsure{(p)}}
\left\{\begin{array}{l}\equiv v_{\varepsilon-1}+1\ (p)\quad(\text{si}\\
\deleted{\unsure{u_\alpha}}\quad N'=\sum v_\alpha p^\alpha)
\end{array}\right.
\end{aligned}
\]

\uline{Dém.} Le développement $p$-adique de $p^\varepsilon-1$ est
\[
p^\varepsilon-1=(p-1)+(p-1)p+\dots+(p-1)p^{\varepsilon-1}
\]
de m
\[
p^{\varepsilon-1}-1=(p-1)+\dots+(p-1)p^{\varepsilon-2}\quad\text{si }\varepsilon\geq2\quad(=0\text{ si }\varepsilon=1)
\]
Prenant
\[
\left.\begin{aligned}
N'+1&=u_0+u_1p+\dots+u_\nu p^\nu\\
N'&=v_0+v_1p+\dots+v_\nu p^\nu
\end{aligned}\right\}\quad(\nu\geq\varepsilon-1)
\]
d'où
\[
\begin{aligned}
N+1&=0+u_0p+u_1p^2+\dots+u_\nu p^{\nu+1}\\
N&=p-1+v_0p+v_1p^2+\dots+v_\nu p^{\nu+1}
\end{aligned}
\]
\uline{!} On trouve
\[
\begin{cases}
c_{N+1,p^\varepsilon-1}\equiv c_{u_0,p-1}\,c_{u_1,p-1}\dots c_{u_{\varepsilon-2},p-1}\quad(p)\quad(\equiv1\ \inserted{(p)}\text{ si }\varepsilon=1)\\
c_{N'+1,p^{\varepsilon-1}-1}\equiv c_{u_0,p-1}\,c_{u_1,p-1}\dots c_{u_{\varepsilon-2},p-1}\quad(p)\quad(\equiv1\ (p)\text{ si }\varepsilon=1)
\end{cases}
\]
\[
\begin{cases}
c_{N,p^\varepsilon}\equiv c_{v_{\varepsilon-1},1}\equiv v_{\varepsilon-1}+1\\
c_{N',p^{\varepsilon-1}}\equiv c_{v_{\varepsilon-1},1}\equiv v_{\varepsilon-1}+1
\end{cases}
\]

\margin{\uline{Rmq} \unsure{$\overline m$|même} si $E$ un $\mathbb F_p$-vectoriel, on \unsure{ne} peut conclure que $x_i=x_N$ $\forall i\in I$ (ex : \inserted{$p=2$,} $N=2$, $e=2$, donc on a $x_2,x_3,x_4,x_6$ on a $x_2=x_3=x_6$ mais \uline{pas} néc. $=x_4$)}
\deleted{\uline{\unsure{Prop}}}\inserted{\uline{Lemme 2}} Soit $\mathbb Z_p$ le localisé de $\mathbb Z$ en $p$, $E$ un $\mathbb Z_p$-module, $N\in\mathbb N$, $e\geq0$ tel que $p^e\geq N+1$, \deleted{\inserted{\unsure{(donc $e\geq1$)}}} $I$ l'ens. des entiers de la forme $N$ et $N+p^\varepsilon$, $0\leq\varepsilon\leq e$, $(x_i)_{i\in I}$ une famille d'éléments de $E$, tel qu'on ait
\figure{\begin{tikzpicture}\draw (0,0)--(0,-2.9);\end{tikzpicture}}
\[
\binom ji(x_j-x_i)=0\qquad i=N,N+1,\quad i<j\in I\quad(\text{donc }j=N+p^\varepsilon,\ 0\leq\varepsilon\leq e)
\]

\mspage{20}
